\subsection{A counterexample to the main implication}
We fix a universe and take all categories to be small in it. We use
covariant functors $H:A\to\Set$, with arrows in $\El_AH$ directed by
$(a,h)\xrightarrow{f}(b,H(f)h)$. We put $BA=|NA|$. A nonempty connected
category $A$ is total if each component of every such category of
elements is weakly contractible. Exactness means exactness of
$\colim_A:\Ab^A\to\Ab$. The module $P_p=\Z[A(p,-)]$ is covariant;
$\Delta\Z$ is the constant right $\Z A$-module. We write
$T_i(M)=\Tor_i^{\Z A}(\Delta\Z,M)=L_i\colim_A M$.

\begin{theorem}[Refuted: totality does not imply exactness, Part A.1]\label{A:main}
Let $\mathcal C_\delta$ have as objects all pairs $([n],E)$, where
$n\geq1$, $[n]=\{1<\cdots<n\}$, $E\subseteq\binom{[n]}3$, and
\[
 \delta_X(U)=2|U|-3|\{e\in E:e\subseteq U\}|>0
 \quad(\varnothing\ne U\subseteq[n]).
\]
An arrow $f:X\to Y$ is a strictly increasing injection which is an
induced hypergraph embedding and satisfies
\[
 \delta_Y(U)>\delta_Y(fX)
 \quad(fX\subsetneq U\subseteq Y).
\]
Identities and composition are those of functions. This is a small,
nonempty connected total category. Let $V$ be the one-point object,
let $W$ be the ordered three-point object with its unique triple as
an edge, and let $v_i:V\to W$ select its $i$th vertex. The three
arrows $v_1,v_2,v_3$ have no integral affine post-equalizer.
Consequently $\colim_{\mathcal C_\delta}:\Ab^{\mathcal C_\delta}\to\Ab$
is not exact.
\end{theorem}
\begin{proof}
The closure calculus in Lemma~\ref{A:predimension-amalgam} proves
that the specified arrows compose and that the category has
amalgamation and common successors. Theorem~\ref{A:predimension-total}
proves totality for every covariant set-valued action, including
infinite actions and actions with noninjective transitions. Its proof
uses the finite alternating-amalgam theorem to classify all open
stabilizers of a rich homogeneous limit, then gives explicit homotopy
comparisons for arbitrary actions. Proposition~\ref{A:predimension-triple}
proves that the displayed triple has no affine equalizer and exhibits
a diagram $Q$ with $T_1(Q)\cong\Z$. These results verify every
hypothesis and the asserted failure for this explicitly specified
category.
\end{proof}
For two ordinary parallel arrows the affine-equalizer assertion
still follows from the audited permutation-module calculation in
Lemma~\ref{B:row-audit}. Thus three is the smallest possible size
of an obstructing ordinary family. The row and column presentations
of Proposition~\ref{B:lead3} now give different answers in the very
same total category; their distinction is not merely a missing
formal step in a proof.

\subsection{An arithmetic reduction}
\begin{proposition}[Proved: fieldwise detection of integral exactness]\label{A:fieldwise}
For every small category $A$, including the empty category and categories
with several components, the following conditions are equivalent:
\begin{enumerate}
\item $\colim_A:\Ab^A\to\Ab$ is exact;
\item $\colim_A:(\Q\text{-}\mathrm{Mod})^A\to\Q\text{-}\mathrm{Mod}$
is exact and, for every prime $p$, the functor
$\colim_A:(\F_p\text{-}\mathrm{Mod})^A\to\F_p\text{-}\mathrm{Mod}$
is exact.
\end{enumerate}
No totality hypothesis is needed. Consequently, Parts A.1 and A.2 taken
together ask equivalent universal questions. A negative answer to A.1
necessarily gives a failure over $\Q$ or over at least one prime field,
although the obstructing diagram need not remain the same after reduction.
\end{proposition}
\begin{proof}
The ordinary abelian-group colimit of a diagram of $R$-modules, for
$R=\Q$ or $\F_p$, carries the unique compatible $R$-module structure,
and its underlying abelian group is the ordinary abelian-group colimit.
This follows from the direct-sum-and-relations construction, since the
subgroup of relations is an $R$-submodule. Forgetting scalars is exact,
so (1) implies (2).

Assume (2). We use the normalized bar complex
\[
 C_n(A;M)=\bigoplus_{a_0\to\cdots\to a_n}^{\mathrm{nondegenerate}}
 M(a_0),
\]
with its usual face differential; its homology is $T_n(M)$. If $M$ is
a $\Q$- or $\F_p$-module diagram, this is, after restriction of scalars,
exactly the bar complex computing derived colimits over that field.
Thus such diagrams have $T_i(M)=0$ for $i>0$. This comparison is an
equality of the specified complexes, not an assumption that arbitrary
base change is flat.

We recall the three closure facts used here. A direct sum or filtered
colimit of acyclic diagrams is acyclic because the bar complex commutes
with these operations and both operations are exact in $\Ab$. If
$0\to U\to V\to W\to0$ is objectwise exact, the bar complexes form a
short exact sequence; their long exact homology sequence shows closure
under extensions and under kernels of epimorphisms between acyclic
diagrams.

Let $M:A\to\Ab$ be arbitrary. For each prime $p$, its $p$-primary
torsion subfunctor is the union of $M[p^n]$, where
$M[p^n](a)=\{m:p^nm=0\}$. The quotient
$M[p^n]/M[p^{n-1}]$ is an $\F_p$-module diagram. Induction on $n$,
then filtered colimits, prove acyclicity of the $p$-primary subfunctor.
The full torsion subfunctor $M_{\mathrm{tor}}$ is their direct sum and
is therefore acyclic. Put $M_{\mathrm{tf}}=M/M_{\mathrm{tor}}$.
It is objectwise torsion free, so the natural map
\[
 M_{\mathrm{tf}}\longrightarrow\Q\otimes_\Z M_{\mathrm{tf}}
\]
is a monomorphism in $\Ab^A$. Its cokernel is objectwise torsion:
each element is a finite rational linear combination, and a common
denominator kills its class. Both the middle diagram and this cokernel
are acyclic. Closure under kernels of epimorphisms makes
$M_{\mathrm{tf}}$ acyclic. Finally the sequence
$0\to M_{\mathrm{tor}}\to M\to M_{\mathrm{tf}}\to0$ makes $M$
acyclic. Therefore $T_1(M)=0$ for every $M$, which is equivalent to
exactness of the right-exact colimit functor.
\end{proof}

\begin{corollary}[Refuted: both coefficient implications, Part A.2]\label{A:coefficients}
The same category $\mathcal C_\delta$ has nonexact colimits over
every nonzero unital coefficient ring $R$. In particular the
implication from totality to exactness is false over $\Q$ and over
$\F_p$ for each prime $p$, with one category and one triple
working simultaneously for all these coefficients. Explicitly, put
\[
 D_R:R[\mathcal C_\delta(W,-)]\longrightarrow
       R[\mathcal C_\delta(V,-)]^2,
 \qquad g\longmapsto(gv_2-gv_1,gv_3-gv_1).
\]
Then $D_R$ is injective, $Q_R=\operatorname{coker}D_R$ has
projective dimension exactly one, and $T_1(Q_R)\cong R$.
\end{corollary}
\begin{proof}
Proposition~\ref{A:predimension-triple} proves injectivity over
every $R$ by a private vertex in each finite nonempty support.
The resulting two-term projective resolution has colimit differential
zero, so $T_1(Q_R)=R\ne0$. Its projective dimension is at most
one by that resolution and is not zero by this nonvanishing.
The totality proof concerns the underlying category and is
independent of $R$.
\end{proof}

\subsection{Proof and construction routes investigated}
The accompanying report records the precise scope of each attempted
route. The model-category observation that totality makes ordinary
simplicial-set colimits homotopy invariant does not by itself identify
the underlying set colimit with the abelian-group colimit. A free
abelian cotriple resolution has permutation modules as its terms, but
some face maps send a basis vector to a linear combination; it is not
the linearization of a simplicial set diagram. We therefore do not infer
acyclicity of its colimit complex from totality. The row/column formulas
below identify the nonzero simultaneous-intersection obstruction in
the separator, and explain why these proof-side routes cannot imply
the proposed universal vanishing.
